% ================================================================
\renewcommand{\realmaccent}{atrAccent}  % reset to default teal
\chapter{For the Game Master}
% ================================================================
% TODO: opening fiction for the GM chapter. No dedicated art asset yet
% either --- every other chapter has a \chapterart image; this one
% doesn't, until one gets generated.
\begin{readaloud}
\lipsum[51][1-3]
\end{readaloud}

\lettrine{R}{unning} \emph{Legend of the Technomancers} means building
worlds, voicing everyone who isn't a Traveller, and deciding what the
dice mean when the outcome is genuinely in doubt. This chapter is yours:
how to prep a realm, how to keep a session moving, and how to adjudicate
the mechanics that make failure matter.

\section{Running the Game}
Keep the world moving, but let the players choose their actions. Describe
what the opposition is doing and what will happen if nobody intervenes;
ask a Traveller what they do about it. You can intrude whenever tension
calls for it --- a guard spots the open door, the ogre bears down on the
bridge, the court begins to turn against a speaker. You need neither an
initiative slot nor a failed player roll to advance the bad guys. Make
the danger clear enough for the players to answer it when they can; if
they ignore it, or the fiction leaves them no chance to stop it, follow
through. The GM never needs to roll for an NPC's attack.

\section{Setting the Target Number}
% REAL RULE, 2026-09-22: relocated from The Shape of a Session, which
% was accumulating too much GM-facing material to stay a section there.
% DRAFT, 2026-09-23 (revised): six-step difficulty ladder and odds
% table. Percentages come from auto_play_testing/dice_odds.csv (targets
% 1--5) plus the same binomial maths for target 6; regenerate them if
% the dice engine changes.
Set the target number from the fiction, before you look at anyone's
sheet. Ask how hard the task is in this situation, pick a number from
the ladder below, and announce it along with the stakes before the
player rolls. The player's Stat and tags already reward them for being
good at it; if you also lowered the target for experts, you would count
their edge twice.

\begin{statblock}[The Difficulty Ladder]
\begin{tabularx}{\linewidth}{@{}l Y@{}}
\dice{1} \keyword{Routine} & An ordinary person succeeds most of the time, but it could still go wrong --- climbing a rough wall, bluffing a bored guard.\\[3pt]
\dice{2} \keyword{Hard} & Needs real skill or good tools --- picking a quality lock, talking down an angry crowd.\\[3pt]
\dice{3} \keyword{Very Hard} & A challenge even for an expert --- tracking prey across bare rock, out-arguing a court wizard in her own hall.\\[3pt]
\dice{4} \keyword{Heroic} & The stuff of stories --- leaping between speeding vehicles, cracking military encryption on the fly.\\[3pt]
\dice{5} \keyword{Legendary} & Only the best, with every edge they can find, have a real hope --- wrestling a dragon to the ground, swaying a god.\\[3pt]
\dice{6} \keyword{Mythic} & Beyond legendary. Even the best will usually fail --- holding back a collapsing mountain, talking the sun out of rising.\\
\end{tabularx}
\end{statblock}

Most rolls in a session should be \dice{1}, \dice{2}, or \dice{3}.
Even at \dice{1} a failure is still possible, and a partial success
still costs something. Save \dice{4} and above for moments you want the
table to remember.

\subsection*{Reading the Odds}
As a rule of thumb, when the pool is \emph{twice the target plus one},
a full success is a coin flip; when it is twice the target minus one,
at least a partial success is a coin flip. A pool smaller than the target cannot
succeed at all --- if you set a target no Traveller can reach, you have
said the action is impossible, so say so plainly instead.

\begin{statblock}[Odds by Target]
Each cell shows the percentage chance of \emph{at least a partial}
success, then of a \emph{full} success or better.
\medskip

\begin{tabularx}{\linewidth}{@{}l Y Y Y Y@{}}
\textbf{Target} & \textbf{3 dice} & \textbf{5 dice} & \textbf{7 dice} & \textbf{9 dice}\\[3pt]
\dice{1} & 88 / 50 & 97 / 81 & 99 / 94 & 99 / 98\\[3pt]
\dice{2} & 50 / 13 & 81 / 50 & 94 / 77 & 98 / 91\\[3pt]
\dice{3} & 13 / 0 & 50 / 19 & 77 / 50 & 91 / 75\\[3pt]
\dice{4} & 0 / 0 & 19 / 3 & 50 / 23 & 75 / 50\\[3pt]
\dice{5} & 0 / 0 & 3 / 0 & 23 / 6 & 50 / 25\\[3pt]
\dice{6} & 0 / 0 & 0 / 0 & 6 / 1 & 25 / 9\\
\end{tabularx}
\end{statblock}

An ordinary person rolls about two dice, which meets a target of
\dice{1} three times in four. A starting Traveller acting in their specialty usually
rolls four to six dice; a Traveller outside their comfort zone might
roll two or three. Even the largest possible pool, \dice{10} dice,
meets a target of \dice{6} less than two times in five.
Critical successes are common at target \dice{1} and rare above
\dice{2}, so low targets are where the spectacular moments live.

\subsection*{Moving the Target}
Circumstances move the target; the character's own edges go in the
pool. Raise the target by one when the situation works against the
Traveller --- darkness, a storm, a ticking clock, a hostile crowd. Lower
it by one, never below \dice{1}, when the situation helps: the party
set up the moment, scouted ahead, or turned the terrain to their
advantage. Rewarding good planning this way is how players learn that
the fiction matters as much as the numbers.

Never change a target after the dice hit the table. If a player
describes a cleverer approach before rolling, re-set the target for
that new approach and announce it again.

\section{Adjudicating Consequences}
% REAL RULE, 2026-09-22: relocated from The Shape of a Session.
When a player fails a roll, you decide what it costs them. The default
is simple: dock one point from the pool matching the Stat they rolled
--- Body for a physical failure, Mind for a mental one, Face for a
social one. If the fiction threatens a different track, use that one
instead; a failed attempt to outwit an armed foe might end in a wound.
Save a bigger hit, two or three points at once, for failures with real
weight behind them: a fall from height, a lie that unravels in front
of the wrong person, a spell that gets away from the caster entirely.

Let the fiction decide the size of the cost, not a fixed table. Ask
yourself what would actually happen here before you touch anyone's
sheet. And remember that pool damage is not the only consequence
available to you --- a failed roll is also licence to complicate the
scene, introduce a new problem, or simply say the attempt didn't work.
Save the pools for failures that leave a lasting mark.

\section{Creative Uses of Abilities}
% DRAFT, 2026-09-25: from a design chat with Chris. Two ideas:
% (1) a successful roll can create a temporary tag for a later roll,
% even one in a different Stat --- blow up the mech, then intimidate its
% friends with an extra die; (2) a table-wide Rule of Cool / Grounded
% dial for stunts, where under Grounded a stunt's side effect joins the
% stakes as benefit on success, harm on failure, both on a partial.
% Created tags are deliberately just tags (+1 die, five-tag cap) so this
% adds one rule, not a family of situational-modifier rules.
Players will use their abilities in ways nobody wrote down: a fireball
aimed at the ground to launch the caster across a chasm, a hacked
elevator dropped forty floors to outrun pursuers. Encourage it. Don't ask
whether the ability's description permits the stunt; ask what resists
it, set a target, and decide what the stunt's own side effects mean.
Refuse a stunt only when the fiction rules it out, and say why, so the
player can try another angle.

\subsection*{Setting Up the Next Roll}
Sometimes the point of an action is not the action itself but the edge
it gives the next one. When a player says so before rolling, a success
creates a new tag: name it, and it adds \dice{1} die to any later roll
it clearly helps, in any Stat and for any character. It works exactly
like a tag on a character sheet, and counts toward the five-tag limit
the same way. A roll in one Stat can therefore feed a roll in another.
\begin{itemize}
  \item A Traveller blows up a mech, then turns to its crew and tells
    them to get lost. Destroying the mech created
    \atrtag{body}{Mech Wrecker}, and it adds a die to the
    \keyword{Face} roll that follows.
  \item A smooth talker needles a duellist until he loses his temper.
    The \keyword{Face} roll leaves him \atrtag{face}{Rattled}, and the
    party's fighter counts that tag when she lunges at him.
  \item A shapeshifter becomes a bear. The transformation earns
    \atrtag{magic}{Bear}, which adds a die whenever a bear's body would
    help: mauling, climbing, roaring down a doorway.
\end{itemize}
The degree of success shapes the tag. A partial success earns the tag
with a cost or complication attached; the shapeshifter who scrapes a
partial might come out as an otter, whose \atrtag{magic}{Otter} helps
in the river but not in the brawl. A failure earns nothing and costs
what failures usually cost. A critical success earns a tag that lasts
longer or that the whole party can use.

A tag lasts as long as the fiction does. A momentary edge --- a
flinching crowd, a rattled duellist --- is spent on the first roll that
uses it. A lasting change, like a bear's body, holds until something
ends it. This overlaps with lowering the target when the party sets up
a moment (see \emph{Moving the Target}): pick one, not both. Lower the
target when the setup happens in passing; ask for a roll and create a
tag when the setup is risky enough to fail.

\subsection*{Set the Dial}
Before play, decide how seriously your table takes side effects, and
tell the players. You can set it once for the campaign or once per realm
--- high fantasy might run loose while cyberpunk runs hard.
\begin{itemize}
  \item \keyword{Rule of Cool}: stunts work the way the player pictures
    them. The fireball launches the caster and doesn't burn them. The
    roll still matters, but only the way any roll does: on a failure the
    stunt doesn't come off and the caster pays the ordinary cost.
  \item \keyword{Grounded}: a stunt carries the consequences its physics
    implies. Before the roll, name the \emph{benefit} the player wants
    and the \emph{harm} the stunt risks. On a success, they get the
    benefit; on a failure, they take the harm; on a partial success,
    they get both. A critical success gives the benefit with no harm,
    plus whatever extra the fiction allows.
\end{itemize}
Neither setting is more correct. Rule of Cool suits pulpy, fast games
where players should feel free to try anything; Grounded suits games
where every trick should cost something, and it makes the players think
twice before the more reckless ones. What matters is that the players
know the setting in advance, so they can plan around it instead of
discovering it after the dice land. When a stunt works once, it should
work the same way the next time; your ruling is now part of the world.

\begin{example}
\emph{The rocket jump.} A wizard needs to cross a chasm, so she fires a
fireball at her own feet. That is two rolls, the same way turning
invisible and then sneaking is. First she makes the blast, a
\keyword{Weird} roll whose benefit is \atrtag{magic}{Launched}. Then
she jumps, a \keyword{Body} roll at \dice{2} with that tag adding a die.
Under Rule of Cool, a failed blast just fizzles, and she can jump
without it or find another way. Under Grounded, the GM names the harm
before the first roll: \dice{1} Body from the blast. On a partial she
is launched but singed; on a failure she is burned and gets no lift.
\end{example}

\begin{example}
\emph{The express elevator.} In a cyberpunk arcology, a hacker overrides
an elevator's brakes to drop the party forty floors ahead of the
security team. The building's system resists, so it's a \keyword{Mind}
roll at \dice{3}. Under Rule of Cool, the brakes catch at the lobby and
everyone walks out. Under Grounded, the harm is the stop at the
bottom: \dice{1} Body to everyone in the car. Note that the harm is
Body even though the roll was Mind. The track the harm hits follows the
fiction, not the Stat the player rolled. On a partial, the party escapes
but staggers out bruised; on a failure, they take the jolt and the car
jams between floors.
\end{example}

\section{Coherence and Catastrophe}
% REAL RULE, 2026-09-22: relocated from The Shape of a Session.
Track the party's \keyword{Coherence} pool openly at the table ---
everyone should watch it drop together, since it's shared. Once a
character's own Stat pool runs dry, route any further loss of that kind
into Coherence instead.

When Coherence hits \dice{0}, something has to give. Treat it as the
payoff for tension that's been building all session, not a random
punishment: the world itself stops holding together around the party.
A reliable go-to is pulling the party loose from the realm entirely and
into the next one --- it closes out a session on a real note of
consequence and hands you a clean opening for the next. You don't have
to use that option every time; anything that reads as \emph{reality
rejecting the party} works. Whatever you pick, make sure the table feels
it as a consequence of what happened, not an arbitrary reset.

\section{Prepping a Realm}
% TODO: how to build a new world before the players arrive.
\lipsum[53][1-4]

\section{Running NPCs}
Give an NPC a \keyword{Body}, \keyword{Mind}, and \keyword{Face} track
when it matters how the players overcome them. A mook has \dice{1}
point in each: one telling blow, one convincing ruse, or one public
humiliation is enough to take them out of the conflict. Do not add the
tracks together --- the players need to bring \emph{one} of them to
\dice{0}, not all three.

Give a monster or other major opponent uneven tracks that express what
makes them dangerous and where they are vulnerable. An ogre might have
Body \dice{4}, Mind \dice{1}, and Face \dice{2}: hard to bring down in
a fight, easy to outwit. Let the players learn that weakness through
what the ogre says and does, rather than treating its numbers as a
secret puzzle.

\subsection*{When an NPC Track Reaches Zero}
Zero means the NPC has lost this conflict \emph{in that way}, not that
all three tracks have been emptied or that every defeat is a death.
\begin{itemize}
  \item \keyword{Body} \dice{0}: they fall down or are otherwise physically
    unable to keep fighting. They might be unconscious, pinned, or
    driven off; death depends on what happened in the fiction.
  \item \keyword{Mind} \dice{0}: they are outwitted. They fall for the
    players' trick, give up the information, or let the players past
    their guard.
  \item \keyword{Face} \dice{0}: they are shamed or lose standing. They
    might withdraw in embarrassment, be laughed at in court, or find
    that nobody will back them any longer.
\end{itemize}
Describe the consequence that fits the action that emptied the track.
An outwitted ogre can still be physically strong; it simply cannot
keep blocking the players with the trick they have already beaten.

\section{Pacing a Session}
% TODO: scenes, turns, and knowing when to cut away.
\lipsum[55][1-4]

\begin{notebox}[Behind the Screen]
% TODO: practical GM-facing checklist --- what to have on hand, prep
% shortcuts, running Coherence and the pools at the table.
\lipsum[56][1-3]
\end{notebox}
